\documentclass[11pt]{article}
\usepackage[margin=1in]{geometry}
\usepackage{mathptmx} % Times New Roman-like font
\usepackage[T1]{fontenc}
\usepackage{enumitem}
\usepackage{parskip}
\usepackage{hyperref}
\hypersetup{colorlinks=false, pdfborder={0 0 0}}

% Section heading: left title (bold) + right-aligned italic subtitle on same line
\newcommand{\sechead}[2]{%
  \vspace{0.35cm}\noindent
  {\large\bfseries #1}\hfill{\large\itshape #2}\par
  \vspace{0.15cm}
}
% Section heading with no subtitle
\newcommand{\secheadonly}[1]{%
  \vspace{0.35cm}\noindent
  {\large\bfseries #1}\par
  \vspace{0.15cm}
}

\setlist[itemize]{leftmargin=0.35in, itemsep=4pt, topsep=4pt}

\pagenumbering{gobble}

\begin{document}

% ---------- Title block ----------
\begin{center}
{\LARGE\bfseries RepairFlow -- Right-to-Repair Assistant \& Reparability Portal}\\[6pt]
Project Proposal for UCS503P\\
Submitted to: Mr. Asif\\
\textbf{Thapar Institute of Engineering and Technology}\\[10pt]
Project Team\\
Punya Chhabra(1024030684) $\vert$ Arnav Gupta(1024030573) $\vert$ Kaushik Jindal (1024030442)\\[10pt]
August 16, 2026
\end{center}

\vspace{0.4cm}

% ---------- 1. Higher-order goal ----------
\sechead{1. Higher-order goal}{}
This project aims to advance the right-to-repair movement by making device repair transparent, accessible, and trustworthy for both consumers and independent repair shops. While the underlying issue of e-waste and repair-ecosystem fragmentation spans an entire industry, the immediate engineering objective is to translate repair-shop discovery, request tracking, and reparability information into a measurable, deployable software solution.

% ---------- 2. Time-to-value ----------
\sechead{2. Time-to-value}{}
\begin{itemize}
  \item We will deliver meaningful increments quickly by prototyping the core repair request and tracking workflow, then validating it with users within a short iteration window.
  \item We will prioritize fast feedback over perfection by using weekly iterations, including CI-backed automation early, to reduce release friction.
  \item Success is defined by what can be delivered and measured in the first iteration: a working repair-shop directory and request-tracking flow. This will then be improved based on user feedback.
\end{itemize}

% ---------- 3. Problem Statement ----------
\secheadonly{3. Problem Statement}{}
Device owners and independent repair shops lack a single, transparent platform to manage repairs, access parts and documentation, and make repair the easier choice. Common pain points include:
\begin{itemize}
  \item No repair visibility: once a device is dropped off, there is no way to track status, cost, or timeline.
  \item Limited parts \& documentation: repair shops struggle to source genuine parts and official repair manuals.
  \item No self-repair option: due to the absence of clear repair instructions, users cannot repair the product themselves and must rely on a third party.
  \item Repair loses to replacement: slow, expensive, and opaque repairs push users toward buying new devices, adding to e-waste.
\end{itemize}
Why this matters now (contextual relevance): Global e-waste hit 62 million tonnes in 2022, with only 22\% properly collected and recycled (UN Global E-waste Monitor). At the same time, right-to-repair laws are gaining ground, from India's Right to Repair Portal to the EU's repair directive and new US state laws, all pushing manufacturers toward repairable design. A simple platform connecting owners, shops, and repair documentation can meaningfully close this gap.

% ---------- 4. Proposed Solution ----------
\sechead{4. Proposed Solution}

\sechead{4.1 Overview}{}
We propose building a web-based ``R2R Assistant'' platform that uses RAG (Retrieval-Augmented Generation) to support right-to-repair, with the following core components:
\begin{itemize}
  \item Verified repair directory: independent repair shops can register and list verified profiles with clear, upfront pricing.
  \item Repair request \& tracking: users can submit a repair request and track its status in real time, from drop-off to completion.
  \item Parts \& documentation hub: a central place to find repair guides, manuals, and genuine spare parts, powered by RAG.
  \item Reparability profiles: each device is scored on how easy it is to repair, based on parts availability and access to guides, helping users make informed decisions.
  \item Sustainability dashboard: shows repairs completed and e-waste avoided, giving real data to support right-to-repair advocacy.
\end{itemize}
\secheadonly{4.2 Core workflow}
\begin{itemize}
  \item User describes a device or fault and the RAG assistant surfaces relevant reparability information, guides, and matching verified shops nearby.
  \item User submits a repair request to a chosen shop and receives a tracking ID.
  \item Shop staff confirm the request, update status and cost through completion, and the record closes out.
  \item The sustainability dashboard updates with the logged repair and estimated e-waste diverted.
\end{itemize}

\secheadonly{4.3 Operational constraints for an educational setting}
\begin{itemize}
  \item Keep the solution usable on low-to-mid range devices via responsive web design.
  \item Avoid dependence on advanced research algorithms beyond a standard RAG pipeline; focus on workflow, data correctness, and retrieval quality.
  \item Design for deployability using common web hosting and managed databases.
\end{itemize}

% ---------- 5. Solution Approach ----------
\sechead{5. Solution Approach}{}
This is an engineering course project, so we focus on scalable data-management and workflow aspects rather than novel algorithm research.

\sechead{5.1 Search \& matching}{}
\begin{itemize}
  \item Reparability and documentation retrieval will use a standard RAG pipeline: relevant repair manuals and guides are retrieved and then an LLM generates the answer, without relying on novel research.
  \item Shop matching and ranking will use simple filtering based on rating, distance, and price from verified shop records.
  \item Repair-status freshness will be tracked using ``last updated'' timestamps on each request, so users can see if a status is stale rather than assuming it's current.
\end{itemize}

\sechead{5.2 Web architecture}{}
A typical 3-tier web architecture:
\begin{itemize}
  \item Frontend: responsive UI for device owners (search, request, track) and repair-shop staff (job and inventory dashboard).
  \item Backend API: authentication, RAG query handling, request/tracking CRUD, and shop verification workflow.
  \item Database: stores users, shops, devices, repair requests, documentation embeddings, and sustainability logs.
\end{itemize}

\secheadonly{5.3 CI/CD and fast delivery enabler}
CI/CD will be used to:
\begin{itemize}
  \item Automatically build and run tests on every change.
  \item Produce repeatable deployment artifacts to staging.
  \item Enable safe and frequent releases of incremental features (e.g., adding the sustainability dashboard after core tracking is stable).
\end{itemize}

% ---------- 6. Evaluation Criterion ----------
\sechead{6. Evaluation Criterion}{}
We will measure success using clear metrics tied to how the system is actually used.

\secheadonly{6.1 Primary evaluation metric}
Time-to-answer (TTA): the time between a user submitting a repair query (device or fault description) and the RAG assistant returning a useful list of guides and nearby verified shops.
\begin{itemize}
  \item Target: median TTA of 3 seconds or less during the pilot.
  \item How we'll measure it: by logging server-side timestamps for each query.
\end{itemize}

\secheadonly{6.2 Secondary evaluation metrics}
\begin{itemize}
  \item Request completion rate: the percentage of repair requests that are completed versus abandoned or cancelled.
  \item Retrieval relevance rate: the percentage of queries that return at least one useful guide or shop, compared to manual search as a baseline.
  \item User clarity: how clear and useful users find the results, based on a simple 1--5 rating.
  \item Reliability: system uptime during business hours, targeting 99\% or higher during the pilot.
\end{itemize}

\sechead{6.3 Pilot validation plan}{}
\begin{itemize}
  \item Run a pilot with 10--20 users and 3--5 partner repair shops (or simulated shop accounts if real partners aren't available).
  \item Compare results before and after the pilot, using short interviews or manual shop calls to establish a baseline.
\end{itemize}

% ---------- 7. Scalability ----------
\sechead{7. Scalability}{}
The proposition should scale in both theoretical and practical senses.

\secheadonly{Scalable (theoretically)}{}
As the number of shops, devices, and users grows, the system is designed to keep up by:
\begin{itemize}
  \item keeping the frontend and backend separate, so each can grow independently,
  \item using pagination and indexing to keep common searches (like device model or shop location) fast,
  \item using a stateless API design and a scalable search index for document retrieval.
\end{itemize}

\secheadonly{Operationally deployable (practically)}{}
The system is built to run on standard, easily available web hosting:
\begin{itemize}
  \item a managed database and vector store,
  \item deployment through containers or a hosting platform,
  \item a CI/CD pipeline to push updates to staging and production.
\end{itemize}

% ---------- 8. Engine availability heuristic ----------
\secheadonly{8. Engine availability heuristic}
A mature set of ``engines'' (tooling/services) is available to implement this as a web-based project:
\begin{itemize}
  \item Web framework for REST APIs and reactive frontend pages.
  \item Managed relational or document database, plus a vector database/embedding store for RAG.
  \item LLM API for retrieval-augmented generation and answer synthesis.
  \item Authentication approach (token-based or session-based) with separate roles for device owners and shop staff.
  \item Test framework for unit/integration tests.
  \item CI runner and staging deployment target.
\end{itemize}
We will use these mature components to focus on workflow design, correctness, and measurement rather than inventing new infrastructure.

% ---------- 9. Project scope and deliverables ----------
\sechead{9. Project scope and deliverables}{}

\sechead{9.1 Initial deliverable}{}
\begin{itemize}
  \item Shop registration and verification form
  \item Repair request submission with RAG-assisted guide/shop suggestions
  \item Owner results page: shop directory with rating, price, distance
  \item Basic logging and timestamps for TTA measurement
  \item CI: build + backend unit tests + linting
  \item CD to staging with a single command or pipeline job
\end{itemize}

\secheadonly{9.2 Subsequent deliverables}
\begin{itemize}
  \item Repair status tracking with notifications (SMS/email optional in later iteration; can start with in-app notifications)
  \item Parts \& documentation hub with genuine spare-part sourcing links
  \item Sustainability dashboard for repairs completed and e-waste diverted
  \item Shop dashboard for job and inventory management
  \item Improved retrieval and indexed queries for scale readiness
\end{itemize}

% ---------- 10. Risks and mitigations ----------
\secheadonly{10. Risks and mitigations}
\begin{itemize}
  \item Data accuracy: shop-submitted status and pricing may become outdated. We'll add ``last updated'' timestamps and send shops periodic reminders to update their listings.
  \item Shop adoption: shops may be slow to adopt the platform if updates feel like a hassle. We'll keep the update process simple, ideally a few taps, and support bulk upload via CSV for parts.
  \item Verification overhead: manually verifying every shop could slow down onboarding. We'll use a lightweight process where shops upload documents and an admin approves them.
  \item Retrieval quality: the RAG assistant may occasionally give inaccurate or made-up answers. We'll rely on curated source documents, show where each answer came from, and let users fall back to contacting the shop directly.
  \item Measurement ambiguity: to avoid confusion when evaluating results, we'll clearly define what counts as each request/tracking state and set up timestamp logging before development starts.
  \item Operational issues during the pilot: we'll test on a staging environment early and make sure CI/CD produces consistent, repeatable deployments.
\end{itemize}

% ---------- 11. Summary ----------
\secheadonly{11. Summary}
This project proposes a deployable RAG-based web platform that gives device owners transparent access to verified repair shops, reparability information, and parts/documentation, while tracking repairs end to end. It meets the course criteria by emphasizing time-to-value through rapid iterations, implementing CI/CD to support frequent safe releases, and using measurable evaluation criteria (especially time-to-answer). It remains focused on engineering workflow, data correctness, and scalability readiness rather than research-driven novelty, while advancing the right-to-repair movement and helping reduce e-waste.

\end{document}
